\documentclass[11pt,a4paper]{article}
\usepackage[margin=27mm]{geometry}
\usepackage{amsmath,amssymb,amsthm,hyperref}
\newtheorem{lemma}{Lemma}
\newtheorem{theorem}[lemma]{Theorem}
\newcommand{\Sz}{\operatorname{Sz}}
\newcommand{\N}{\operatorname{N}}
\newcommand{\C}{\operatorname{C}}
\title{A rank obstruction in prime-degree Suzuki groups}
\author{Human Compression project}\date{}
\begin{document}\maketitle
\begin{abstract}
We prove the nonexistence assertion in Theorem 5.4 of the fixed 1999
rank-three Suzuki paper. The argument uses its subgroup definition of
residual weak primitivity directly, so assumes neither flag transitivity
nor a geometric interpretation of every residue. Two successive maximal
intersections suffice. The required proper-overgroup information is taken
from the classical Suzuki subgroup classification, with the elementary
point-stabilizer calculation imported from the separately reviewed note.
\end{abstract}

\section{The subgroup condition}
Let $I$ be finite and let $(G_i)_{i\in I}$ be subgroups of a finite group
$G$. Write $G_J=\bigcap_{j\in J}G_j$ and $G_\varnothing=G$.
The condition used in the source \cite{source}, printed page~46, is
\begin{equation}\label{rw}
 \text{for every }J\subsetneq I\text{ there is }i\in I\setminus J
 \text{ such that }G_{J\cup\{i\}}\text{ is maximal proper in }G_J.
\end{equation}
Here $J=\varnothing$ is included. The native statement of Theorem~5.4,
printed page~55, excludes rank \emph{at least} three, not merely rank
strictly greater than three.

\begin{lemma}\label{strict}
Under \eqref{rw}, for every $J\subsetneq I$ and every $i\notin J$,
$G_{J\cup\{i\}}<G_J$. In particular all $G_i$ are proper, no $G_i$
is contained in a different $G_j$, and adjoining any further type to an
intersection strictly decreases it.
\end{lemma}
\begin{proof}
If $G_J\le G_i$, then $G_{I\setminus\{i\}}\le G_J\le G_i$.
Consequently $G_{I\setminus\{i\}}=G_I$, contradicting \eqref{rw}
at $I\setminus\{i\}$, where the only available type is $i$.
This proves the asserted strictness and its consequences.
\end{proof}

\section{Classical inputs and the point stabilizer}
Put $q=2^n$, with $n\ge3$ odd, and $G=\Sz(q)$. We retain the
classical Suzuki subgroup and ovoid structure used in the source and
specified in the accepted notes \cite{preliminary,normalizers}.
The maximal subgroups have types
\[
 B=P\rtimes C_{q-1},\quad D=C_{q-1}\rtimes C_2,\quad
 N_\epsilon=C_{a_\epsilon}\rtimes C_4,\quad
 \Sz(2^d),
 \qquad a_\epsilon=q+\epsilon2^{(n+1)/2}+1,
\]
where the last type has $d\ge3$ and $n/d$ prime. Every proper
subgroup is contained in one of these groups. The three odd torus
orders are pairwise coprime. Every nonidentity odd element is in a
unique cyclic torus; its centralizer is that torus. A split torus has
normalizer $D$, whose involution inverts it. In $N_\epsilon$ the
order-four complement acts without nonidentity fixed points on its
odd kernel, and its square acts by inversion. The group $P$ has
order $q^2$ and centre $Z$ elementary abelian of order $q$. A nonidentity
involution fixes exactly one ovoid point; $B$ is a point stabilizer.
These are attributed structural premises, not new classification claims.

The reviewed elementary proof in \cite{preliminary} shows that every
maximal subgroup of $B$ contains $Z$. We record its consequence for
\eqref{rw}, separately from any flag-transitivity hypothesis.

\begin{lemma}\label{borel}
If $|I|\ge2$ and \eqref{rw} holds, no $G_i$ is a point stabilizer $B$.
\end{lemma}
\begin{proof}
Suppose $G_0=B$. Choose $i\ne0$ such that $J=B\cap G_i$ is
maximal in $B$. Then $Z\le J$. No dihedral or nonsplit torus
normalizer contains $Z$: their Sylow two-subgroups have order two
or are cyclic of order four. Nor does a proper simple subfield group
$\Sz(s)$ contain $Z$, since $q\ge s^3>s^2$, the order of its Sylow
two-subgroups. A point stabilizer containing $Z$ must equal $B$,
because any fixed nonidentity involution in $Z$ has a unique fixed
ovoid point. Thus $B$ is the only maximal subgroup of $G$ containing
$J$. The proper subgroup $G_i$ lies in a maximal subgroup, so $G_i\le B$,
contrary to Lemma~\ref{strict}.
\end{proof}

\section{Two maximal intersections}
For the remainder assume $n$ is prime. There are then no proper simple
subfield groups. This hypothesis is essential to the following
proper-overgroup arguments.

\begin{lemma}\label{overgroups}
If $x$ has nontrivial odd order dividing $a_\epsilon$, every proper
subgroup containing $x$ lies in the unique $N_\epsilon$ associated
with its torus. If a dihedral subgroup $E$ has a nontrivial odd rotation
subgroup in a split torus, every proper subgroup containing $E$ lies in
the normalizer $D$ of that split torus.
\end{lemma}
\begin{proof}
For the first assertion, take a maximal subgroup containing the given
proper subgroup. Coprimality excludes point stabilizers, split
normalizers and the opposite nonsplit type. In a nonsplit normalizer
an odd-order element belongs to its odd kernel. Uniqueness of the torus
containing $x$ therefore identifies that maximal subgroup.

For the second assertion, coprimality excludes both nonsplit types.
A point stabilizer cannot contain $E$. Indeed, its quotient by its
normal two-subgroup is abelian of odd order. An odd rotation has
nontrivial image in this quotient, whereas the involution inverting it
has trivial image. The inversion relation would force that odd image
to equal its inverse, a contradiction. The remaining maximal subgroup
is a split normalizer, and the rotation subgroup determines its unique
torus, hence determines this normalizer.
\end{proof}

\begin{lemma}\label{nonsplit}
If $n$ is prime and $|I|\ge3$, no $G_i$ is a full nonsplit torus
normalizer $N=C_a\rtimes C_4$.
\end{lemma}
\begin{proof}
Suppose $G_0=N$. For $i\ne0$, the subgroup $K_i=G_0\cap G_i$
cannot contain a nonidentity odd-order element: Lemma~\ref{overgroups}
would give $G_i\le N$, violating Lemma~\ref{strict}. Thus $K_i$ is
a two-group, cyclic of order at most four. Choose $i$ with $K_i$
maximal in $N$. It cannot have order one or two, since it lies in a
Sylow two-subgroup of order four. Hence $K_i\cong C_4$.

Apply \eqref{rw} at $\{0,i\}$ to choose $k\notin\{0,i\}$ such
that $G_{0ik}$ is maximal in $K_i$; it has order two. An involution
$z\in N$ lies in a unique cyclic subgroup of order four. To see this,
conjugate its Sylow subgroup to the specified complement $V$. Then
$z$ is its involution and $\C_N(z)=V$, because $z$ inverts the odd
kernel. Any cyclic group of order four containing $z$ centralizes $z$
and therefore equals $V$.
It follows that the cyclic two-subgroup $K_k$ containing $G_{0ik}$
is contained in $K_i$. Hence $G_{0ik}=G_{0k}$, again contradicting
Lemma~\ref{strict}.
\end{proof}

\begin{lemma}\label{split}
Suppose $n$ is prime and $q-1$ is composite. If $|I|\ge3$, no $G_i$
is the full split normalizer $D=C_{q-1}\rtimes C_2$.
\end{lemma}
\begin{proof}
Suppose $G_0=D$, and write $A=C_{q-1}$. Each $K_i=G_0\cap G_i$,
$i\ne0$, is either contained in $A$ or is a subgroup of order at most
two. Indeed, any other subgroup of the odd dihedral group $D$ is
dihedral with a nontrivial odd rotation subgroup; Lemma~\ref{overgroups}
would force $G_i\le D$.
Choose $i$ such that $K_i$ is maximal in $D$. The maximal subgroups
of $D$ are $A$ and dihedral groups of order $2(q-1)/p$, where $p$
is a prime divisor of $q-1$. The latter have nontrivial odd rotation
subgroups because $q-1$ is composite, so $K_i=A$.
Now choose $k\notin\{0,i\}$ with $G_{0ik}$ maximal in $A$.
Its order is $(q-1)/p>1$ for some prime $p$, so is odd and nontrivial.
The preceding restriction on $K_k$ forces $K_k\le A=K_i$. Thus
$G_{0ik}=G_{0k}$, contrary to Lemma~\ref{strict}.
\end{proof}

\begin{theorem}\label{main}
Let $n\ge3$ be an odd prime and suppose $2^n-1$ is composite.
There is no subgroup family for $G=\Sz(2^n)$ satisfying \eqref{rw}
with $|I|\ge3$. In particular there is no residually weakly primitive
pre-geometry of rank at least three in the source's sense.
\end{theorem}
\begin{proof}
The case $J=\varnothing$ in \eqref{rw} makes at least one $G_i$
maximal in $G$. The classical list leaves $B,D,N_+,N_-$, since
$n$ is prime. Lemmas~\ref{borel}, \ref{nonsplit} and \ref{split}
exclude all four possibilities.
\end{proof}

This proves exactly the source's Theorem 5.4, including its full rank
range. The point-stabilizer exclusion also proves the subgroup form of
its quoted Lemma 4.1. The nonsplit and split arguments above are only
claimed in prime field degree; they do not by themselves dispose of
all composite-degree assertions in Lemmas 4.2--4.6. No sufficient
realization has been used as a necessity statement. No finite group
search, new conjecture resolution, Lean verification or compression
measurement is claimed for this component.

\begin{thebibliography}{9}
\bibitem{source}D.~Leemans, \emph{The rank 3 geometries of the simple
Suzuki groups $\Sz(q)$}, Note di Matematica 19(1) (1999), 43--63.
Fixed twenty-one-page edition, especially printed pp.~46 and~55.
\bibitem{preliminary}Human Compression project, \emph{Maximal subgroups
and elementary reductions for Suzuki geometries}, accepted preliminary
component for L115, with its explicit classical Suzuki premises.
\bibitem{normalizers}Human Compression project, \emph{Normalizer fibres
in rank-two Suzuki geometries}, accepted L115 component. Its structural
inputs are attributed there to Suzuki, Tits and L\"uneburg.
\end{thebibliography}
\end{document}
